\subsection{Strategic Approach}  
The proposed solution addresses the critical security vulnerabilities of the Baltic region by fundamentally shifting the paradigm of early warning capabilities from a reactive, ground-based radar posture to a proactive, space-based observation framework. Ground-based systems are inherently limited by the curvature of the Earth, terrain masking, and fixed geographical positioning, which severely restricts their line-of-sight against low-altitude or highly compressed ballistic trajectories. This limitation manifests as the "insufficient warning times" currently challenging Baltic national security authorities. 

To overcome these limitations, the concept relies on a hybrid satellite constellation composed of 6 \abbr{GEO} satellites and 72 \abbr{LEO} satellites deployed at an altitude of approximately 1,200 km and evenly distributed across 8 orbital planes with 9 satellites per plane as the current Space Segment concept.\footnote{The Space Segment team reports having performed \abbr{GMAT} and \abbr{VTS} constellation analyses. Detailed model files, assumptions, coverage plots and quantitative evidence are to be incorporated or referenced when delivered. Continuous global observation or guaranteed 60-second warning performance should not be treated as independently verified in this report until that evidence has been reviewed.} This specific architecture was selected over an alternative \abbr{MEO}/\abbr{LEO} framework based on a rigorous trade study. While \abbr{MEO} layers offer wider coverage outside the equatorial region, operating within the inner Van Allen radiation belt severely compromises sensor lifespans, drives up lifecycle costs and introduces critical detection gaps during initial launch sequences. The chosen \abbr{GEO}/\abbr{LEO} hybrid mitigates these risks. The \abbr{GEO} segment is primarily dedicated to detecting the initial booster phase of a missile launch, when the thermal signature is strongest and earliest detection is possible. The \abbr{LEO} segment then resolves polar coverage limitations and provides follow-on observation, enabling continuous tracking of the missile trajectory and supporting a more accurate preliminary impact-area estimation. 

Rather than waiting for a ballistic missile to clear international borders or rise to a high altitude to be detected by ground radars, this concept focuses entirely on capturing the massive electromagnetic and thermal energy emitted during the initial ignition sequence (the boost phase). By detecting launch events as early as possible and maintaining uninterrupted tracking thereafter, the system is designed to provide national authorities with timely warning information for crisis management, civil protection, and emergency response. In addition, the Ground Segment securely archives mission data for algorithm improvement, calibration, performance verification, operator training, and post-event technical analysis, consistent with R-004. 

\subsection{System Architecture}
\label{sec:ground-segment-boundary} 
The system baseline consists of two integrated segments: the Space Segment and the Ground Segment. The Ground Segment includes mission command and control, the Mission Control Center, ground stations, data processing, the Early Warning System and data archiving as specified by the existing WBS. Authorized governmental authorities are external recipients of warning products, not a separate User and Data Segment. The first warning is generated automatically only after data from the two selected \abbr{LEO} satellites have been evaluated on the ground. 


\subsection{The Space Segment}
The Space Segment forms the core sensing layer of the system and consists of the proposed 6 \abbr{GEO} satellites and 72 \abbr{LEO} satellites. The \abbr{GEO} satellites provide persistent wide-area observation and are optimized for the earliest detection of missile launches during the boost phase, when the thermal plume is most intense and the launch location can be identified rapidly. When a \abbr{GEO} sensor detects a candidate thermal burst, it performs preliminary onboard denoising and transmits the event location and supporting detection information to the Ground Segment. Ground Segment command-and-control functions identify two suitable \abbr{LEO} satellites according to event position, observation geometry, spacecraft readiness and connectivity. Commands are delivered using available ground-station paths or \abbr{ISL} relay connections when supported. The selected \abbr{LEO} satellites undertake follow-on observation, may communicate with other satellites when necessary, and return encrypted observation data and a preliminary trajectory estimate to the Ground Segment. 

The combined constellation is intended to support timely detection of events that might threaten the Baltic region, including events originating outside that region; quantitative detection and coverage performance is to be confirmed using the Space Segment simulation evidence. The sensors aboard the satellites are designed to detect transient infrared signatures associated with missile launches and to support automated cueing of follow-on observations. 


\subsection{The Ground Segment} 
The Ground Segment coordinates observations and ground-directed \abbr{LEO} selection, transmits tracking commands, receives secure spacecraft telemetry, processes observations, and sends warning information to authorized national authorities. It includes a Ground Station Network, Command and Control System, Mission Control Center, Data Processing System, Early Warning System and Data Archiving System. Satellite operations and the authority-facing alert interface are therefore supported without a separate User and Data Segment. 

Secure ground stations receive the initial \abbr{GEO} report and subsequent \abbr{LEO} data and forward validated packets to the Data Processing System. Computationally intensive multi-satellite fusion and the refined trajectory/impact-area estimation are performed by the central Ground Segment processing infrastructure; edge preprocessing at individual stations is optional where technically justified. Once observations from two selected \abbr{LEO} satellites have been assessed, the Early Warning System automatically generates and dispatches a preliminary, uncertainty-labelled warning to authorized authorities. Operator confirmation, correction or cancellation follows the initial dispatch; the Mission Control Center records this review. The Data Archiving System securely retains mission records for the purposes permitted by R-004. A 60-second delivery target is retained but requires a verified end-to-end latency budget. 


\subsection{Ground Segment Warning Dissemination and Data Archiving}
The Early Warning System implements automated communication of the preliminary, LEO-informed warning product to authorized national crisis-management and civil-protection authorities, including later confirmed, updated or cancelled versions. National authority interfaces are the final EWAS delivery boundary; downstream public alert activation remains outside EWAS. A separate User and Data Segment is not retained in the project baseline.

The Data Archiving System, also inside the Ground Segment, manages raw and processed mission data, event and warning records, and relevant audit logs. Access and use of archived mission data are restricted to algorithm improvement, calibration, performance verification, operator training and post-event technical analysis under R-004. More general strategic exploitation or sharing of historical archives is not an assumed EWAS function.

\subsection{Core Functional Pillars}  
The mission requirements are met through three distinct, conceptually separated functional operations within the system architecture. 

\paragraph{Pillar I: Boost-Phase Infrared Detection}
The system's earliest warning capability is based on infrared detection of the missile's boost phase. The primary technical mechanism for achieving the "earliest possible" detection relies on \abbr{IR} sensing. During launch, ballistic missiles generate a strong thermal plume from burning propellants, producing a distinct infrared signature that can be detected from space. The \abbr{GEO} satellites are primarily responsible for this initial detection function, providing persistent observation while executing preliminary denoising and reporting the candidate event to the Ground Segment for rapid selection and cueing of the tracking layer. 

\paragraph{Pillar II: Trajectory Tracking and Impact Estimation}
Once a launch is detected, the \abbr{LEO} constellation takes over the tracking function. The concept utilizes multi-satellite observation. After receiving the initial \abbr{GEO} report, the Ground Segment selects two suitable \abbr{LEO} satellites and cues them through the best available supported ground or \abbr{ISL}-relay path. The selected satellites then orient their sensors, acquire complementary observations, estimate a preliminary trajectory and return encrypted data to the ground. Using multiple observation angles the satellites provide the geometric data needed to estimate the missile's trajectory with increasing confidence.  

\paragraph{Pillar III: Low-Latency Dissemination and Sovereign Archiving}
Once the Ground Segment has processed the observations of two selected \abbr{LEO} satellites and derived an impact-area estimate, the Early Warning System automatically transmits a clearly marked preliminary warning to an authorized national authority interface. It does not wait for prior operator confirmation. The Mission Control Center subsequently supports confirmation, updates and cancellation. The Data Archiving System preserves the relevant sensor, processing and warning records for the limited technical purposes specified in R-004. The end-to-end objective of 60 seconds from first \abbr{GEO} detection to delivered preliminary warning remains an unverified design driver pending latency analysis. 